\appendix

\section{Canonical 数据合同的诊断方法}\label{app:canonical-diagnostic-chain}

本附录按误差来源汇总 canonical 数据合同的诊断步骤及其物理解释。

\begin{table}[!ht]
\centering
\caption{Canonical 数据合同的异常诊断与误差来源隔离。}
\label{tab:canonical-diagnostic-chain}
\scriptsize
\begin{tabularx}{\textwidth}{@{}l>{\raggedright\arraybackslash}X>{\raggedright\arraybackslash}X@{}}
\toprule
诊断步骤 & 观察与分析 & 科学解释 \\
\midrule
同轨迹比较 & 修正跨轨迹比较后，mode1 满足判据而 mixed 仍未满足。 & 消除比较对象不一致造成的假阳性。 \\
重复单模态工况 & 单一 mode1 重复也可能出现闭合误差。 & 模态混合不是异常的必要条件。 \\
余切投影 & \(N=2\) 截断不足并非根因，projected momentum 的闭合更好。 & 动量应按固定基投影到余切空间。 \\
求解器输出语义 & generalized-alpha 内部 stage 不可由公共 API 直接取得；BDF \texttt{tsteps} 可输出同一 solver step 的状态。 & BDF 仅用于状态采样和标签生成。 \\
固定基初值 & 对位移和速度应用 normalization/sign transform 后，\(12/12\) 初值一致且 mixed 短窗闭合。 & 异常来自初值与固定基定义不一致。 \\
双层时间步比较 & G2 从 \(0.00147337\) 增至 \(0.00155680\)。 & 两层结果不足以满足预设单调收敛判据。 \\
三层与 \(1\,\mathrm{s}\) 验证 & 观察阶为 \(1.64186\)--\(2.00088\)，延长工况的闭合与能量指标满足阈值。 & 数据合同满足后续多幅值实验的数据生成要求。 \\
\bottomrule
\end{tabularx}
\end{table}


诊断过程依次检查比较对象、模态混合、降阶截断、余切动量定义、求解器输出语义和固定基初值。跨轨迹比较产生的假阳性说明，闭合误差必须在同一物理轨迹和同一时间采样上计算。单模态重复工况仍可能出现异常，因而不能把误差直接归因于模态混合。增大降阶表示后，异常并未表现出由 \(N=2\) 截断主导的特征；相比 \(M_N\dot q_{\mathrm{native}}\)，投影动量给出更好的闭合。

generalized-alpha 的内部 stage 不能通过公共 API 直接取得，因此不能保证导出的位移、速度和载荷来自同一内部求解阶段。二阶 BDF 的 \texttt{tsteps} 输出用于获得同一 solver step 的状态和向量场标签。该选择不意味着 BDF 是辛积分器，也不把离散轨迹解释为辛映射。

力闭合与扰动能量漂移采用彼此独立的判据。早期力比较对端点和动量锚点敏感；在公共支持 indices \(2\ldots998\) 和固定 \(p[2]\) 锚点下，三条轨迹的 force-v2 联合 L2 均低于 \(2\%\)。这一结果不改变原 \(0.5\%\) Bridge 扰动能量判据的逐轨迹结论。

\section{指标定义与数值设置}\label{app:numerical-settings}

\subsection{Canonical 映射与闭合量}

固定 \(N=2\) 基下采用
\begin{align}
q &= \Pi_{\mathrm{conf}}(U), &
p &= B^{T}P_{\mathrm{full}}, \\
\dot q &= M_N^{-1}p, &
\dot p &= B^{T}L .
\end{align}
其中 \(q\) 属于构形坐标，\(p\) 属于对应的余切坐标。闭合误差分别比较导出的 \(\dot q\) 与 \(M_N^{-1}p\)、导出的 \(\dot p\) 与 \(B^{T}L\)，并同时考察 \(q\)、\(p\) 和 total energy 的时间离散收敛。

三层时间步计算采用 coarse/middle/fine 网格。若相邻两层误差分别为 \(e_h\) 和 \(e_{h/2}\)，观察阶定义为
\[
r_{\mathrm{obs}}=\log_2\!\left(\frac{e_h}{e_{h/2}}\right).
\]
所得观察阶为 \(1.64186\)--\(2.00088\)。在 \(1\,\mathrm{s}\) 工况上，G2、G3 点态和 G3 impulse 的阈值均为 \(0.02\)，扰动能量峰峰漂移阈值为 \(0.005\)。

\subsection{Bridge 与低维网络实验}

Bridge 扰动能量漂移按轨迹分别与 \(0.5\%\) 判据和探索性的 \(1\%\) 判据比较；force-v2 使用独立的 \(2\%\) 联合 L2 判据。两类指标不合并为单一结论。

P0 使用四维非线性耦合振子，训练/验证/测试轨迹数为 \(80/16/20\)，每条 80 步，\(\Delta t=0.025\)，训练 300 epochs。五个随机种子用于估计算法重复性。

P1 和 P1.5 各包含 \(60/60\) 次完成的训练与评价运行，P1.6 包含 \(30/30\) 个硬约束模型。三条物理轨迹 mode1、mode2 和 mixed12 按 leave-one-whole-trajectory-out 划分；不同随机种子共享同一物理轨迹集合，因此不能用于扩大物理样本量。P1.5 比较 physics-only 与 residual 学习，P1.6 比较平衡硬约束前后的精度和结构残差。

\subsection{后续数值实例化}

多幅值统计、完整 96 轨迹数据集和 P2 模型比较将沿用整轨迹划分，并报告跨工况均值、方差和配对比较。开放端口研究还需给出功率平衡、耗散符号和端口泄漏指标；高保真研究需量化 Kirchhoff 模型相对 Timoshenko 梁和实体模型的偏差。数值 KAM 分析则需在统一范数下给出非退化裕度、Diophantine 裕度、高阶误差、辛映射缺陷和不变环面残差。实现代码和补充数据将在配套代码与补充材料中提供。
